\ExplSyntaxOn \makeatletter

\keys_define:nn { hduthesis / exam }
  {
    watermark .tl_set:N = \l__hduthesis_exam_watermark_tl,
  }
\NewDocumentCommand \hdutexset {m} { \keys_set:nn { hduthesis / exam } {#1} }
\DeclareHookRule { shipout / background }
  { hduthesis / exam } { before } { pgfrcs }
\AddToHook { shipout / background } [ hduthesis / exam ]
  {
    \tl_if_empty:NF \l__hduthesis_exam_watermark_tl
      {
        \tikz [ remember~picture, overlay ]
          \node [ opacity = .2 ] at (current~page)
            {
              \includegraphics [ width = \linewidth ]
                { \l__hduthesis_exam_watermark_tl }
            };
      }
  }

\RequirePackage
  {
    geometry, mathtools, array, bm, esint, keytheorems, libertine, inconsolata,
    siunitx, fixdif, derivative, physics2, tasks, anyfontsize
  }
\RequirePackage[inline]{enumitem}
\RequirePackage{fontawesome5, fancyqr, circuitikz}
\hook_gput_code:nnn { cmd/faIcon/before } { . }
  {
    \group_begin:
    \fontsize { \fp_eval:n { .9 * \f@size } } { 0 } \selectfont
  }
\hook_gput_code:nnn { cmd/faIcon/after } { . } { \group_end: }
\geometry{margin = 1in}

\RequirePackage[svgnames]{xcolor}
\newkeytheorem{note}
  [ style = remark, numbered = false,
    name = \color{DarkRed}\sffamily \faIcon{info-circle}~Tip. ]
\newkeytheorem{problem}
  [ style = definition,
    name = \color{DarkGreen}\sffamily \faIcon{pen-square}~Problem ]
\newkeytheorem{solution}
  [ style = definition, numbered = false,
    name = \color{MidnightBlue}\sffamily \faIcon{check-square}~Solution ]
\usephysicsmodule{ab, op.legacy}
\usephysicsmodule{bm-um.legacy}
\settasks { label = \sffamily\Alph*., item-indent = 2.5em, label-offset = .5pt }
\newcommand \true {\addtocounter{task}{1}\faIcon{check-circle}}
\def \@result#1 {\underline{~#1~}}
\def \s@result {\textsf{【\qquad 】}}
\DeclareRobustCommand \result {\@ifstar\s@result\@result}
\keys_define:nn { textfig }
  {
    ratio      .tl_set:N    = \l__textfig_ratio_tl,
      ratio    .initial:n   = 7 \c_colon_str 2,
    position   .tl_set:N    = \l__textfig_position_tl,
      position .initial:n   = c,
    unknown    .code:n      = \tl_if_novalue:nF {#1}
      { \tl_set_eq:NN \l__textfig_position_tl \l_keys_key_tl }
  }
\NewDocumentEnvironment{textfig}{ O{} }
  {
    \group_begin:
    \keys_set:nn { textfig } {#1}
    \seq_set_split:Nee \l__whats_textfig_ratio_seq { \c_colon_str }
      { \l__textfig_ratio_tl }
    \noindent
    \minipage [ \l__textfig_position_tl ]
      {
        \fp_eval:n
          {
            .96*\seq_item:Nn \l__whats_textfig_ratio_seq {1}/
              (
                \seq_item:Nn \l__whats_textfig_ratio_seq {1} + 
                \seq_item:Nn \l__whats_textfig_ratio_seq {2}
              )
          }\dimexpr \linewidth
      }
  }
  {
    \endminipage
    \vspace{4.5pt plus 1.5pt minus 1.5pt} \par
    \group_end:
  }
\NewDocumentCommand \textfigsplit {}
  {
    \endminipage
    \hspace{.02\linewidth}
    \minipage [ \l__textfig_position_tl ]
      {
        \fp_eval:n
          {
            .96*\seq_item:Nn \l__whats_textfig_ratio_seq {2}/
              (
                \seq_item:Nn \l__whats_textfig_ratio_seq {1} + 
                \seq_item:Nn \l__whats_textfig_ratio_seq {2}
              )
          }\dimexpr \linewidth
      }
  }

% \hook_gput_code:nnn { cmd/maketitle/before } { . } { \vskip -\baselineskip }
% \hook_gput_code:nnn { cmd/maketitle/after } { . }
%   { \__weyoung_maketitle_aftercode_aux: }
\cs_new_protected:Nn \__weyoung_maketitle_aftercode_aux:
  {
    \vspace*{-2\baselineskip}
    \begin{center}
      \begin{minipage}{.66\textwidth}
        \centering
        \begin{tasks}
          [ label-width = 5\ccwd, item-indent = 2cm,
            after-item-skip = 0pt ](2)
          \task[\textbf{课程信息}：] 理学院/A0715022
          \task[\textbf{授课教师}：] 陈许敏
          \task[\textbf{解析提供}：] \textsf{Skyrmion}
          \task[{\makebox[5em][s]{\textbf{\textsf\TeX~排版}：}}] \textsf{Skyrmion}
        \end{tasks}
      \end{minipage}
      \begin{minipage}{.15\textwidth}
        \centering
        \begin{tikzpicture}
          \node [ inner~sep = 0 ] (wechat) {
            \fancyqr
              [ size = 4em, image~padding = 1,
                image = \raisebox{4em/33}{\textcolor{DarkGreen}{\faIcon{weixin}}},
                r~color = DarkGreen, version = 4
              ]{http://weixin.qq.com/r/sRJiemTEGjllrSAZ90ej}
            };
          \node [ inner~sep = .5ex, align = center, right, font = \scriptsize ]
            at (wechat.east) {未\\央\\学\\社};
        \end{tikzpicture}
      \end{minipage}
      \begin{minipage}{.15\textwidth}
        \centering
        \begin{tikzpicture}
          \node [ inner~sep = 0 ] (qq) {
            \fancyqr
              [ size = 4em, image~x~padding = 2, image~y~padding = 1,
                image = \raisebox{4em/33}{\textcolor{MidnightBlue}{\faIcon{qq}}},
                l~color = MidnightBlue, version = 4
              ]{https://qm.qq.com/q/PH7CsQWvm2}
            };
          \node [ inner~sep = .5ex, align = center, left, font = \scriptsize ]
            at (qq.west) {物\\理\\营};
        \end{tikzpicture}
      \end{minipage}
    \end{center}
  }

\newcommand \iu {{\mathup i\mkern1mu}}
\newcommand \eu {\mathup e}
\AtBeginDocument{%
  \let \Re \relax \let \Im \relax
  \DeclareMathOperator{\Re}{Re}
  \DeclareMathOperator{\Im}{Im}
}

\makeatother \ExplSyntaxOff
\usetikzlibrary{ fadings, patterns, arrows, decorations.markings }
\tikzset
  {
    > = stealth',
    every picture/.append style =
      { line join = round, line cap = round, thick }
  }